
\documentclass[a4paper]{article}

\usepackage{url}
\usepackage{amsmath}
\usepackage{verbatim}   		% Useful for program listings
\usepackage[T1]{fontenc}       	% For Swedish characters ÅÄÖ etc.
\usepackage[utf8]{inputenc}
% \usepackage[swedish]{babel} % For Swedish hyphenation
\usepackage{fancyvrb}           	% For lists with tabulators
\fvset{tabsize=4}              	 	% Tabulator size
\fvset{fontsize=\small}         	% List font size
\usepackage{graphicx}		% Imports the graphicx package, useful for images

% generate clickable references & toc
\usepackage[]{hyperref}
\hypersetup{
%    pdftitle={Your title here},
%    pdfauthor={Your name here},
%    pdfsubject={Your subject here},
%    pdfkeywords={keyword1, keyword2},
    bookmarksnumbered=true,
    bookmarksopen=true,
    bookmarksopenlevel=1,
    colorlinks=true,
    pdfstartview=Fit,
    pdfpagemode=UseOutlines,    % this is the option you were lookin for
    pdfpagelayout=TwoPageRight
}

% define namedlabel as per http://texblog.org/2012/03/21/cross-referencing-list-items/
% this is used for naming the requirements.
\usepackage{enumitem, hyperref}
\usepackage{nameref}


%% Macros for mapping
% - use case to subsection level
% - scenario to subsubsection level
% - requirement to paragraph level
%
% Using these results in it being possible to cross-reference
% requirements and getting unique identifiers which are author-chosen.

% usage: \usecase{cross-reference}{display-reference} text text ...
\newcommand{\usecase}[2]{
  \subsection*{#2}
  \label{#1}
  \addcontentsline{toc}{subsection}{\nameref{#1}}
}


% usage: \scenario{cross-reference}{display-reference} text text ...
\newcommand{\scenario}[2]{
  \subsubsection*{#2}
  \label{#1}
  \addcontentsline{toc}{subsubsection}{\nameref{#1}}
}

 % usage: \requirement{cross-reference}{display-reference} text text ...
\newcommand{\requirement}[2]{
  \paragraph*{#2:}
  \label{#1}
  \addcontentsline{toc}{paragraph}{\nameref{#1}}
}



\title{Dragonfly Quadrotor UAV \\ Product Requirements Specification}
\author{Daniel Stenberg \\ Nina Khayyami \\ Daniel Nilsson \\ Eduardo Riffo \\ Adam Steineck}

\date{\today}
%\date{December, 2015}         		% Today's date if not specified


\begin{document}                	% Start of document

\maketitle                      	% Prints the title defined above with \title, \author and \date

\begin{center}
\vspace{64pt}
\includegraphics[scale=1.6]{images/AF_Logotype20141_Black.png}
\vspace{16pt}
\\ \large ÅF Technology
\end{center}

\newpage

\tableofcontents				% Insert table of contents

\newpage

\section{Introduction}

The \emph{Dragonfly} project is an internal competence enhancement project for ÅF Technology employees. The goal is to combine technology, competence and experience from various engineering fields in order to construct a highly advanced quadrotor UAV system. This document contains all product requirements tied to a scenarios belonging to a certain use case.

\section{Use cases}


\usecase{uc:navigate}{UC:Navigate}
\scenario{sc:cxn.sunny}{SC:Nav-Sunny} GPS and map data available.
\requirement{req:nav.sunny.1}{Nav.Sunny.1} THe DF shall navigate along a route.
\requirement{req:nav.sunny.2}{Nav.Sunny.2} The DF shall navigate to start position.

\scenario{sc:cxn.rainy}{SC:Nav-Rainy} GPS not available.
\requirement{req:nav.rainy.1}{Nav.Rainy.1} User shall be notified that navigation is unavailable with cause.

\usecase{uc:cxn}{UC:Connection}

\scenario{sc:cxn-ok}{SC:Cxn-Sunny}
\requirement{req:cxn.sunny.wifi}{Cxn.Sunny.1}: The device should have a functional WiFi connection.
\requirement{req:can.sunny.wless}{Cxn.Sunny.2}: The device should have a functional wireless connection.


\usecase{uc:flight}{UC:Flight}
\scenario{sc:flight-sunny}{SC:Flight-RC}
\requirement{req:PRS-CF001}{PRS-CF001} The user must be able to control the Dragonfly manually through a RC transmitter.
\requirement{req:PRS-CF002}{PRS-CF002} The manual control through the RC must have the highest priority (over mobile app and navigation control).
\requirement{req:PRS-CF003}{PRS-CF003} The user must be able to switch between controlling through mobile app or RC and vice verse while flying.
\requirement{req:PRS-CF004}{PRS-CF004}: When the user releases the RC control levers the Dragonfly must return to a hovering position.

\scenario{sc:flight-sunny}{SC:Flight-App}
\requirement{req:PRS-CF005}{PRS-CF005} The user must be able to control the Dragonfly manually through a mobile app.
\requirement{req:PRS-CF006}{PRS-CF006} The app must be able to detect when the user releases the control as per \nameref{req:PRS-CF007}.

\scenario{sc:flight-sunny}{SC:Flight-All}
\requirement{req:PRS-CF007}{PRS-CF007} The Dragonfly must have a self stabilizing control algorithm.
\requirement{req:PRS-CF008}{PRS-CF008} The user must be able to control the Dragonfly in directions:
\begin{itemize}
\item Forwards
\item Backwards
\item Left
\item Right
\item Rotate clockwise \& anticlockwise viewed from above.
\item Up
\item Down
\end{itemize}
\requirement{req:PRS-CF009}{PRS-CF009}: The Dragonfly must be able to land on its own.
\requirement{req:PRS-CF010}{PRS-CF010}: The Dragonfly must be able to fly in winds up to 10 m/s. % is this average or gust wind? The latter is more realistic as 10m/s is already a 5 on the Beaufort scale.


\section{Sub-system: Charging}

\scenario{sc:Quadcopter unit initiated initiated  charging}{SC:Quadcopter unit initiated charging}
\requirement{req:PRS-CH001}{PRS-CH001}: The Quadcopter unit must have receiver coil attached to receiver cuircuit and a balanced server charger.
\requirement{req:PRS-CH002}{PRS-CH002}: The Quadcopter receiver and transmitter must be placed within a fixed distance.
\requirement{req:PRS-CH003}{PRS-CH003}: The Quadcopter unit must implement a software based capacity monitor.
\requirement{req:PRS-CH004}{PRS-CH004}: The charging base must consist of a main power supply, transmitter circuit and transmitter coil.
\requirement{req:PRS-CH005}{PRS-CH005}: Charging base must enable power supply manually User interaction is mandatory for this process, since it involves tweaking the voltage.
\scenario{sc:Quadcopter user initiated charging}{SC:Quadcopter user initiated charging}
\requirement{req:PRS-CH006}{PRS-CH006}: Enable balanced server charger manually.
\requirement{req:PRS-CH007}{PRS-CH007}: Investigate if Ichanger used in project can be controlled automatically  through software.
\requirement{req:PRS-CH008}{PRS-CH008}: Motion sensors needs to be enabled to guaranty effective charging.
\requirement{req:PRS-CH009}{PRS-CH009}: Mechanism to turn the power supply off should be implement if possible. Monitoring of charging process must be done to avoid overheating and unnecessary  power consumption.
\requirement{req:PRS-CH010}{PRS-CH010}: Charging notifications must be implemented and sent to users or quadcopter unit.
\scenario{sc:Suggest alternative charging base}{SC:Suggest alternative charging base}
\requirement{req:PRS-CH011}{PRS-CH011}: Suggest the nearest charging base to user, but let the user make decision.
\scenario{sc:Redirect to alternative charging base}{SC:Redirect to alternative charging base}
\requirement{req:PRS-CH012}{PRS-CH012}: Automatic charging redirection to charging base must be supported by Quadcopter unit.
\requirement{req:PRS-CH013}{PRS-CH013}: Charging base should be covered by plastic glass with a predetermined thickness.

\section{Sub-system: Safety}

\scenario{sc:safety.detection}{SC: Quadcopter unit detects a human or an object blocking the route; the motor speed is lowered and the quadcopter unit enters hoovering-mode.}
\requirement{req:PRS-SA001}{PRS-SA001} The quadcopter unit shall have a detection unit.
\requirement{req:PRS-SA002}{PRS-SA002} The quadcopter unit shall detect humans and objects that appear within a specified distance from the quadcopter unit.
\requirement{req:PRS-SA003}{PRS-SA003} The quadcopter unit shall have a hoovering-mode, where:
\begin{itemize}
\item motor speed is lowered.
\item
\end{itemize}
\scenario{sc:safety.avoidance}{SC: Quadcopter unit detects a human or an object blocking the route; the route is changed to avoid the obstacle}
\requirement{req:PRS-SA002}{PRS-SA002} Bla bla bla
\scenario{sc:safety.panicmode}{SC: If the safety subsystem fails to avoid humans or objects, the quadcopter automatically switches to panic mode, where motors are stopped and the latest position is made remotely available.}
\requirement{req:PRS-SA001}{PRS-SA001} The quadcopter unit must have a panic-mode, where:
\begin{itemize}
\item engines are switched off.
\item last known position of the device is made available remotely.
\end{itemize}



\section{Example}
\subsection{It's still possible to use the vanilla sections}
And they will appear like use cases etc in the table of contents.
\usecase{uc:cross-reference}{UC:display-reference} Use case text. Use case text.Use case text.Use case text.Use case text.
\scenario{sc:cross-reference}{SC:display-reference} Scenario text. Scenario text. Scenario text. Scenario text. Scenario text. Scenario text.
\requirement{req:cross-reference}{REQ:display-reference} Requirement text. Requirement text. Requirement text. Requirement text. Requirement text. Requirement text.


\section{Example: Referring to requirements}

\subsection{Referring to newcommand uc, scenario and requirements.}
It is preferred to use the \verb!\nameref! to refer to the use cases, scenarios and requirements.
Thus, we look at use case \nameref{uc:navigate} scenario \nameref{sc:cxn.sunny} and requirements \nameref{req:cxn.sunny.wifi}.

\subsection{Sandbox of referring to newcommands}
\nameref{uc:cross-reference}, scenario \nameref{sc:cross-reference} and reqirement \nameref{req:cross-reference}.


\begin{thebibliography}{99}
\bibitem{stenberg} Model-based Design Development and Control of a Wind Resistant Multirotor UAV, C. Månsson, D. Stenberg, Lunds Tekniska Högskola 2014
\end{thebibliography}

\end{document}                  % End of document
